\section{Related Work}
\label{sec:related_work}

Our research intersects four prominent domains of NLP literature: (1) hallucination definition and benchmarking in LLMs, (2) multilingual evaluation and low-resource factual disparities, (3) script transliteration and subword tokenization bottlenecks, and (4) Retrieval-Augmented Generation (RAG) distractor sensitivity and selective abstention.

\subsection{Hallucination Benchmarking in LLMs}
Factual hallucination has long been identified as a critical vulnerability in natural language generation~\citep{maynez2020faithfulness,ji2023survey}. In foundational LLM literature, hallucinations are categorized into \emph{intrinsic} hallucinations (contradicting provided source context) and \emph{extrinsic} hallucinations (asserting ungrounded claims that cannot be verified or falsified by the context). Early benchmarks such as TruthfulQA~\citep{lin2022truthfulqa} evaluated mimicry of human falsehoods and superstitions. More recent efforts have emphasized atomic granularity: FActScore~\citep{min2023factscore} decomposes biographies into atomic propositions evaluated against Wikipedia, while HaluEval~\citep{li2023halueval} synthesizes large-scale QA, dialogue, and summarization errors using sampling-and-filtering pipelines. Furthermore, FAVA~\citep{mishra2024fava} introduces fine-grained span-level error taxonomy and retrieval-augmented correction, while SummaC~\citep{laban2022summac} standardizes NLI-based inconsistency detection in abstractive summarization. However, nearly all existing benchmarks are constructed in English or assume high-resource Western cultural priors, neglecting the syntactic and semantic failure modes prevalent in lower-resource South Asian languages.

\subsection{Multilingual Evaluation and Regional Disparities}
Recent benchmarking initiatives have begun to illuminate severe performance disparities between English and non-English generations. \citet{ahuja2023mega} introduced MEGA, benchmarking generative AI across 16 languages and highlighting pronounced accuracy deficits in low-resource settings. \citet{bang2023multitask} presented a multi-task and multilingual audit of ChatGPT, showing that hallucination rates escalate drastically on low-resource and non-Latin prompts. Poly-FEVER~\citep{zhu2024polyfever} extended claim-verification datasets across 11 languages to study cross-lingual factuality. While foundational multilingual benchmarks such as TyDi QA~\citep{clark2020tydi} and FLORES-101~\citep{goyal2022flores} pioneered typologically diverse evaluation, they focus predominantly on extractive span QA or machine translation rather than generative hallucination and epistemic abstention. Furthermore, existing multilingual datasets frequently rely on translated English trivia, omitting localized administrative, legal, and public-health realities specific to countries like Bangladesh. In such specialized domains, cross-lingual English-pivot retrieval strategies (e.g., translating queries to English, retrieving from global corpora, and back-translating) fail fundamentally: sovereign legal codes (such as the Bangladesh Labour Act or constitutional amendments) and domestic public health circulars have zero representation in English web corpora, rendering translation-based retrieval pipelines ineffective. Matibhrom addresses this gap by curating culturally and institutionally authentic Bengali benchmarks constructed directly from sovereign constitutional statutes, national labor codes, and indigenous clinical guidelines.

\subsection{Script Transliteration and Tokenizer Fragmentation}
In languages employing complex Brahmic-derived scripts such as Bengali, standard subword tokenization constitutes a major structural bottleneck. \citet{hasan2023banglatokenizer} demonstrated that common tokenizers (e.g., Llama and Mistral SentencePiece models) fragment Bengali morphemes into disjoint byte tokens, driving severe token inflation and degrading downstream representations. In parallel, digital South Asian discourse is characterized by widespread informal Latin transliteration, commonly termed \emph{Banglish}. Recent studies~\citep{saha2024robustness,chowdhury20245dialects} demonstrate that LLMs experience significant performance degradation when processing transliterated dialectal inputs. However, previous investigations have evaluated transliteration strictly on classification tasks (e.g., sentiment analysis or machine translation). To our knowledge, Matibhrom is the first benchmark to evaluate the interaction between Latin transliteration, dense multilingual retriever recall, and generative hallucination under grounded retrieval.

\subsection{RAG Vulnerability and Selective Abstention}
Retrieval-Augmented Generation (RAG)~\citep{lewis2020rag} is widely deployed as a practical mitigation against parametric hallucination. However, recent findings emphasize that retrieval is not a panacea: \citet{cuconasu2024power} showed that introducing semantically adjacent or noisy distractor passages frequently misleads LLMs into generating false conclusions. Most recently, \citet{park2024robust} demonstrated that imperfect retrieval---encompassing unanswerable, adversarial, and conflicting documents---severely degrades RAG systems, frequently precipitating catastrophic factual hallucination. Similarly, mtRAG~\citep{katsis2025mtrag} established that handling unanswerability under multi-turn dialogue is a primary bottleneck in current RAG pipelines. When confronted with irrelevant or absent evidence, trustworthy systems must possess the capability to \emph{selectively abstain}~\citep{kadavath2022language,cole2023selectively,yin2023large}. While selective prediction has been formalized for classification and closed-book English QA, its efficacy under multilingual RAG with explicit negative constraints remains largely unexamined. Matibhrom directly fills this void by quantifying how model scale governs the emergence of selective abstention under noisy retrieval in low-resource environments.
